\documentclass[a4paper,12pt]{article}

\usepackage[T2A]{fontenc}
\usepackage[utf8]{inputenc}
\usepackage[russian]{babel}
\usepackage{amsmath,amssymb,mathtools}
\usepackage{helvet}
\renewcommand{\familydefault}{\sfdefault}
\usepackage{icomma}
\usepackage{geometry}
\geometry{left=18mm,right=18mm,top=18mm,bottom=20mm}
\usepackage{microtype}
\usepackage{tikz}
\usetikzlibrary{calc,arrows.meta,positioning,shapes.geometric}
\usepackage{tabularx,booktabs,longtable}
\usepackage{enumitem}
\usepackage{hyperref}
\usepackage{parskip}
\usepackage{fancyhdr}
\usepackage{setspace}
\usepackage{xcolor}

\definecolor{accent}{HTML}{245B65}
\definecolor{darkaccent}{HTML}{153B42}
\definecolor{textgray}{HTML}{2E3336}
\definecolor{softgray}{HTML}{EEF2F2}
\definecolor{warn}{HTML}{8B4A3D}

\hypersetup{colorlinks=true,linkcolor=darkaccent,urlcolor=darkaccent}
\onehalfspacing
\setlength{\parindent}{0pt}
\setlist[itemize]{leftmargin=1.6em,itemsep=0.35em,topsep=0.35em}
\setlist[enumerate]{leftmargin=1.8em,itemsep=0.55em,topsep=0.4em}
\newcommand{\checkitem}{\item[$\square$]}
\newcommand{\key}[1]{\textcolor{darkaccent}{\textbf{#1}}}
\newcommand{\formula}[1]{\[\displaystyle #1\]}

\pagestyle{fancy}
\setlength{\headheight}{14pt}
\fancyhf{}
\fancyhead[L]{\small Степени в задачах прикладного характера}
\fancyhead[R]{\small 30.09.2026}
\fancyfoot[C]{\small \thepage}
\renewcommand{\headrulewidth}{0.4pt}
\renewcommand{\footrulewidth}{0pt}

\tikzset{
  startstop/.style={ellipse,draw=accent,line width=0.9pt,align=center,minimum width=35mm,minimum height=10mm,fill=softgray},
  process/.style={rounded corners=2mm,draw=accent,line width=0.9pt,align=center,text width=90mm,minimum height=12mm},
  decision/.style={diamond,aspect=2.3,draw=accent,line width=0.9pt,align=center,text width=43mm,inner sep=1.2mm},
  arrow/.style={-{Latex[length=3mm]},line width=0.9pt,draw=darkaccent}
}

\begin{document}

\begin{titlepage}
\thispagestyle{empty}
\centering
\vspace*{26mm}
{\Large\textcolor{accent}{Профильная математика · ЕГЭ}\par}
\vspace{11mm}
{\fontsize{25}{30}\selectfont\bfseries Чек-лист\par}
\vspace{5mm}
{\fontsize{20}{25}\selectfont\bfseries Степени в задачах прикладного характера\par}
\vspace{18mm}
{\large Даниил Кичук\par}
\vspace{5mm}
{\normalsize 30 сентября 2026 года\par}
\vfill
{\small Лёвин Артём Александрович эксклюзивно для Даниила Кичука\par}
\vspace{12mm}
\begin{tikzpicture}[x=1cm,y=1cm]
  \draw[accent,line width=1.1pt]
    (-7,0) .. controls (-4.8,1.15) and (-2.4,-1.05) .. (0,0)
           .. controls (2.2,1.05) and (4.8,-1.0) .. (7,0);
\end{tikzpicture}
\end{titlepage}

\section*{1. Главная идея занятия}

В прикладных задачах формула обычно уже дана. Основная работа состоит из четырёх действий: \key{перевести условие на язык математики}, \key{аккуратно подставить данные}, \key{свести вычисления к удобным степеням} и \key{извлечь искомую величину}. На занятии этот маршрут применялся к степенным зависимостям, показательным уравнениям и формулам с третьей и четвёртой степенью.

\begin{itemize}
\checkitem Я отделяю данные от вопроса и выписываю только величины, которые реально входят в формулу.
\checkitem Я заранее отмечаю, какая величина неизвестна и в какой степени она стоит.
\checkitem Если встречаются десятичные коэффициенты при степенях десяти, я при необходимости переписываю их в более удобном виде.
\checkitem Составные числа стараюсь узнавать как степени простых оснований: $625=5^4$, $3125=5^5$, $128=2^7$, $125=5^3$.
\checkitem Если искомая величина находится в степени, применяю обратную степень ко \emph{всей} соответствующей части равенства.
\checkitem Если в условии есть «не менее», «не более», «максимально», сначала проверяю монотонность зависимости и только затем использую граничное равенство.
\end{itemize}

\section*{2. Степени: что должно включаться автоматически}

\subsection*{2.1. Степень степени}
\formula{\left(a^m\right)^n=a^{mn}}
Именно это позволяет убрать показатель у неизвестной. Например, если при $V>0$
\formula{V^{5/3}=A,}
то возводим обе части в степень $\frac35$:
\formula{V=\left(A\right)^{3/5}.}

\subsection*{2.2. Один и тот же множитель в степенном виде}
Если получается число вроде $3125$, полезно разложить его на простые множители:
\formula{3125=5\cdot625=5^5.}
После этого выражения вида $\left(5^5\right)^{3/5}$ считаются без громоздких корней.

\subsection*{2.3. Степени десяти}
Перенос десятичной запятой компенсируется изменением показателя степени десяти:
\formula{6{,}25\cdot10^5=625\cdot10^3,\qquad 5{,}7\cdot10^{-8}=57\cdot10^{-9}.}
Это особенно удобно, когда коэффициенты затем сокращаются.

\section*{3. Типовые модели, разобранные на занятии}

\subsection*{3.1. Степенная зависимость $pV^k=\mathrm{const}$}
На доске рассматривался случай
\formula{pV^{5/3}=200,\qquad p\ge 6{,}25\cdot10^5.}
При $V>0$ величина $V$ уменьшается при росте $p$, поэтому максимальный объём достигается при минимально допустимом давлении:
\formula{p=6{,}25\cdot10^5=625\cdot10^3.}
Тогда
\formula{V^{5/3}=\frac{200}{625\cdot10^3}=\frac1{5^5},\qquad
V=\left(\frac1{5^5}\right)^{3/5}=\frac1{5^3}=0{,}008.}

\textbf{Контроль:} равенство на границе используется потому, что зависимость монотонна. Слова «не менее» и «не более» нельзя механически заменять знаком «равно» в любой задаче.

\subsection*{3.2. Показательный закон убывания}
Для модели
\formula{m(t)=m_0\cdot2^{-t/T}}
на доске использовались $m_0=40$, $T=10$, $m=5$:
\formula{40\cdot2^{-t/10}=5\quad\Longrightarrow\quad
2^{-t/10}=\frac18=2^{-3}\quad\Longrightarrow\quad t=30.}

\textbf{Ключевой приём:} привести обе части к одному основанию и приравнять показатели.

\subsection*{3.3. Два состояния одной системы}
Если для процесса выполняется
\formula{pV^a=\mathrm{const},}
то для двух состояний:
\formula{p_1V_1^a=p_2V_2^a.}
На занятии объём уменьшался вдвое, а давление увеличивалось в четыре раза:
\formula{V_2=\frac{V_1}{2}\;\Longleftrightarrow\;V_1=2V_2,\qquad p_2=4p_1.}
Подстановка даёт
\formula{p_1(2V_2)^a=(4p_1)V_2^a\quad\Longrightarrow\quad2^a=4=2^2,}
откуда $a=2$.

\textbf{Особенно важно:} фразы «уменьшилось в $n$ раз» и «увеличилось в $n$ раз» сначала переводятся в формулы с индексами $1$ и $2$.

\subsection*{3.4. Четвёртая степень в формуле излучения}
На доске использовалась формула
\formula{P=\sigma ST^4,}
где
\formula{\sigma=5{,}7\cdot10^{-8},\qquad S=\frac{10^{20}}{128},\qquad P=1{,}14\cdot10^{25}.}
Удобная подготовка данных:
\formula{\sigma=57\cdot10^{-9},\qquad P=114\cdot10^{23}.}
После подстановки и сокращений:
\formula{T^4=2^8\cdot10^{12}.}
Значит,
\formula{T=\left(2^8\cdot10^{12}\right)^{1/4}=2^2\cdot10^3=4000.}

\subsection*{3.5. Кубическая зависимость}
В задаче о силе натяжения использовалась формула
\formula{T=\rho g l^3-mg.}
При $\rho=1000$, $g=10$, $m=25$ и ограничении $T\le1000$ максимальная допустимая длина достигается на границе $T=1000$:
\formula{1000=1000\cdot10\cdot l^3-25\cdot10,}
\formula{l^3=\frac{1250}{10000}=\frac{125}{1000}=\left(\frac5{10}\right)^3,\qquad l=0{,}5\text{ м}.}

\section*{4. Частые ошибки}

\begin{itemize}
\checkitem \textbf{Потерять скобки при подстановке.} Если $V_1=2V_2$, то $V_1^a=(2V_2)^a$, а не $2V_2^a$.
\checkitem \textbf{Применить степень только к одному множителю.} $(ab)^n=a^nb^n$.
\checkitem \textbf{«Домножить показатель» только слева.} Возведение в степень — операция над обеими частями равенства.
\checkitem \textbf{Считать десятичные числа слишком рано.} Сначала ищите удобную степенную структуру и сокращения.
\checkitem \textbf{Игнорировать направление ограничения.} Для «максимального» или «минимального» значения нужно понимать, растёт или убывает искомая величина.
\checkitem \textbf{Забыть единицы в ответе.} Численные вычисления можно вести без повторения единиц в каждой строке, если исходные единицы согласованы; итоговый ответ записывается с требуемой единицей.
\end{itemize}

\clearpage
\thispagestyle{fancy}
\section*{Алгоритм: задача прикладного характера со степенями}
\vspace{8mm}
\begin{center}
\begin{tikzpicture}[node distance=10mm]
  \node[startstop] (start) {Прочитать условие};
  \node[process,below=of start] (data) {Выписать формулу, данные, ограничение и искомую величину};
  \node[process,below=of data] (translate) {Перевести текстовые отношения в математику: «в $n$ раз», индексы $1/2$, границы};
  \node[process,below=of translate] (prepare) {Подготовить числа: степени десяти, разложение составных чисел, удобные основания};
  \node[process,below=of prepare] (sub) {Подставить данные в скобках и выразить степень искомой величины};
  \node[decision,below=12mm of sub] (exp) {Искомая величина осталась в степени?};
  \node[process,below=13mm of exp] (root) {Если да: возвести обе части в обратную степень};
  \node[process,below=of root] (check) {Проверить знак, ограничение, разумность результата и единицы};
  \node[startstop,below=of check] (finish) {Записать ответ};

  \draw[arrow] (start) -- (data);
  \draw[arrow] (data) -- (translate);
  \draw[arrow] (translate) -- (prepare);
  \draw[arrow] (prepare) -- (sub);
  \draw[arrow] (sub) -- (exp);
  \draw[arrow] (exp) -- node[right]{\small да} (root);
  \draw[arrow] (root) -- (check);
  \draw[arrow] (exp.east) -- ++(24mm,0) node[midway,above]{\small нет} |- (check.east);
  \draw[arrow] (check) -- (finish);
\end{tikzpicture}
\end{center}
\clearpage

\section*{5. Тренировка — Путь А}
\textit{10 заданий. Решения выполняйте отдельно; в этом блоке оставлены только условия.}

\begin{enumerate}
\item Масса вещества меняется по закону
\[
m(t)=64\cdot2^{-t/6}.
\]
Через сколько единиц времени масса станет равна $8$?

\item Для некоторого состояния выполняется
\[
pV^{3/2}=54.
\]
Найдите $V$, если $p=2$.

\item В процессе выполняется $pV^a=\mathrm{const}$. Объём газа уменьшили в $3$ раза, а давление увеличилось в $9$ раз. Найдите $a$.

\item Мощность излучения вычисляется по формуле
\[
P=\sigma ST^4.
\]
Найдите $T$, если $\sigma=5\cdot10^{-8}$, $S=2\cdot10^6$, $P=1{,}6\cdot10^{12}$.

\item Сила натяжения определяется формулой
\[
T=\rho g l^3-mg.
\]
При $\rho=800$, $g=10$, $m=20$ допустимо $T\le800$. Найдите максимально допустимое $l$ в метрах.

\item Масса вещества меняется по закону
\[
m(t)=160\cdot2^{-t/4}.
\]
Найдите момент времени, когда масса станет равна $10$.

\item Для некоторой модели
\[
pV^{5/3}=972.
\]
Найдите $V$, если $p=4$.

\item В процессе выполняется $pV^a=\mathrm{const}$. Объём газа уменьшили в $4$ раза, а давление увеличилось в $8$ раз. Найдите $a$.

\item Для излучающего тела используется формула
\[
P=\sigma ST^4.
\]
Найдите $T$, если $\sigma=8\cdot10^{-8}$, $S=2\cdot10^6$, $P=1{,}296\cdot10^{13}$.

\item Сила натяжения определяется формулой
\[
T=\rho g l^3-mg.
\]
При $\rho=1000$, $g=10$, $m=20$ допустимо $T\le1960$. Найдите максимально допустимое $l$ в метрах.
\end{enumerate}

\clearpage
\section*{Ответы}
\renewcommand{\arraystretch}{1.35}
\begin{center}
\begin{tabularx}{0.78\linewidth}{@{}cXcX@{}}
\toprule
\textbf{№} & \textbf{Ответ} & \textbf{№} & \textbf{Ответ}\\
\midrule
1 & $18$ & 6 & $16$\\
2 & $9$ & 7 & $27$\\
3 & $2$ & 8 & $\dfrac32$\\
4 & $2000$ & 9 & $3000$\\
5 & $0{,}5$ м & 10 & $0{,}6$ м\\
\bottomrule
\end{tabularx}
\end{center}

\vfill
\begin{center}
\small\textcolor{textgray}{Лёвин Артём Александрович эксклюзивно для Даниила Кичука}
\end{center}

\end{document}
